%! Parametric Functions — Unit 8, Lesson 8.1 — Guided Notes (Answer Key)
\documentclass[10pt]{article}
\usepackage{precalculus-article}
\usepackage{precalculus-key}   % pulls in -boxes; do NOT also load -boxes
\newcommand{\TallMath}[1]{$\displaystyle #1\rule[-1.4em]{0pt}{3.2em}$}

\begin{document}

\pageheader{Unit 8, Lesson 8.1}{Guided Notes --- Answer Key}
\namedateperiod

\vspace{0.10in}

%----------------------------------------------------------------------
% Objectives
%----------------------------------------------------------------------
\begin{objectivebox}
\textbf{I will be able to\ldots}
\begin{enumerate}[leftmargin=1.5em, itemsep=4pt]
  \item \textbf{LO 4.1.A} --- Construct a table of values and a graph
    for a parametric function $f(t) = (x(t),\, y(t))$ by evaluating both
    component functions at selected values of the parameter $t$, plotting
    the resulting ordered pairs in order of increasing $t$, and
    interpreting domain restrictions as start and end points.
    \ans{Build table column-by-column; plot in $t$-order; endpoints at $t = a$ and $t = b$.}
\end{enumerate}
\end{objectivebox}

\vspace{0.08in}

%----------------------------------------------------------------------
% Vocabulary
%----------------------------------------------------------------------
\begin{vocabbox}
\begin{description}[leftmargin=0pt, itemsep=6pt]
  \item[\termblanklong{parameter}]
    The \ans{independent} variable $t$ in a parametric function;
    it controls both $x$ and $y$ at the same time.

  \item[\termblanklong{parametric function}]
    A function $f(t) = (x(t),\, y(t))$ in $\mathbb{R}^2$ consisting of
    \ans{two} component equations, both written in terms of a single
    variable $t$.

  \item[\termblanklong{parametric curve}]
    The set of ordered pairs $(x(t),\, y(t))$ traced in the $xy$-plane
    as $t$ \ans{increases} through its domain.
    The curve has a \ans{direction (orientation)} determined by
    increasing $t$.

  \item[\termblanklong{start / end point}]
    When the domain is $a \le t \le b$, the curve begins at
    \ans{$(x(a),\, y(a))$} and ends at \ans{$(x(b),\, y(b))$}.
\end{description}
\end{vocabbox}

\vspace{0.08in}

%----------------------------------------------------------------------
% Hook
%----------------------------------------------------------------------
\begin{hookbox}
\textbf{Entry Question:} A drone's position at time $t$ seconds is
$x(t) = t^2 - 4$ (meters east of launch pad) and $y(t) = 2t$ (meters
north) for $0 \le t \le 4$.

\vspace{4pt}
\textbf{Q1.} What are the drone's coordinates when $t = 0$? When $t = 2$?

\ans{At $t = 0$: $(0^2-4,\,2\cdot0) = (-4,\,0)$. At $t = 2$: $(4-4,\,4) = (0,\,4)$.}

\textbf{Q2.} Could you describe the drone's path with a single equation
$y = f(x)$? What information would that equation \emph{not} tell you?

\ans{Yes: from $x = t^2-4$ we get $t = \sqrt{x+4}$, so $y = 2\sqrt{x+4}$. But the equation loses \emph{time} --- it doesn't tell you where the drone is at any particular moment.}
\end{hookbox}

\vspace{0.08in}

%----------------------------------------------------------------------
% LO 4.1.A — Table and Graph
%----------------------------------------------------------------------
\begin{notesbox}{LO 4.1.A --- Parametric Functions: Table and Graph}

\textbf{Definition:} A parametric function $f(t) = (x(t),\, y(t))$
consists of \ans{two} component equations. Each value of the parameter
$t$ produces one ordered pair $(x(t),\, y(t))$ in the $xy$-plane.

\vspace{8pt}
\textbf{Example:} Let $f(t) = (t + 1,\; t^2 - 3)$ for
$-2 \le t \le 2$.

\vspace{6pt}
\textbf{Step 1 --- Build the table.}

\vspace{4pt}
\begin{center}
\renewcommand{\arraystretch}{1.6}
\begin{tabular}{|c|c|c|c|}
\hline
$t$ & $x(t) = t + 1$ & $y(t) = t^2 - 3$ & $(x,\; y)$ \\
\hline
$-2$ & \ans{$-1$} & \ans{$1$}  & \ans{$(-1,\;1)$}  \\
\hline
$-1$ & \ans{$0$}  & \ans{$-2$} & \ans{$(0,\;-2)$}  \\
\hline
$0$  & \ans{$1$}  & \ans{$-3$} & \ans{$(1,\;-3)$}  \\
\hline
$1$  & \ans{$2$}  & \ans{$-2$} & \ans{$(2,\;-2)$}  \\
\hline
$2$  & \ans{$3$}  & \ans{$1$}  & \ans{$(3,\;1)$}   \\
\hline
\end{tabular}
\end{center}

\vspace{8pt}
\textbf{Step 2 --- Plot and connect.} Plot each $(x, y)$ ordered pair.
Connect in order of increasing $t$; arrows show direction.

\vspace{6pt}
\begin{center}
\begin{tikzpicture}[scale=0.78]
  \draw[linegray, step=1cm] (-2,-4) grid (4,2);
  \draw[->] (-2.3,0) -- (4.5,0) node[right] {\small $x$};
  \draw[->] (0,-4.3) -- (0,2.5) node[above] {\small $y$};
  \foreach \x in {-2,-1,1,2,3,4} {
    \draw (\x, 0.06) -- (\x, -0.06);
    \node[below, font=\scriptsize] at (\x, 0) {$\x$};
  }
  \foreach \y in {-4,-3,-2,-1,1,2} {
    \draw (0.06, \y) -- (-0.06, \y);
    \node[left, font=\scriptsize] at (0, \y) {$\y$};
  }
  % Parametric curve: f(t)=(t+1, t^2-3) for t in [-2,2]
  % Points: (-1,1),(0,-2),(1,-3),(2,-2),(3,1)
  \draw[keyred, very thick, -latex]
    plot[smooth, tension=0.6] coordinates
    {(-1,1) (0,-2) (1,-3) (2,-2) (3,1)};
  % Mark points
  \foreach \pt in {(-1,1),(0,-2),(1,-3),(2,-2),(3,1)} {
    \fill[keyred] \pt circle (2.5pt);
  }
  % Start label
  \node[above left, font=\scriptsize, keyred] at (-1,1) {start};
  % End label
  \node[above right, font=\scriptsize, keyred] at (3,1) {end};
\end{tikzpicture}
\end{center}

\vspace{4pt}
The start point is \ans{$(-1,\,1)$} (at $t = -2$) and the
end point is \ans{$(3,\,1)$} (at $t = 2$).

\end{notesbox}

\vspace{0.08in}

%----------------------------------------------------------------------
% LO 4.1.A — Solving for t
%----------------------------------------------------------------------
\begin{notesbox}{LO 4.1.A --- Solving for the Parameter (Skill 1.A)}

\textbf{Key idea:} To find \emph{when} the curve reaches a specific
$x$- or $y$-value, set the component equation equal to that value and
solve for $t$.

\vspace{8pt}
\textbf{Example (using $f(t) = (t+1,\; t^2-3)$, $-2 \le t \le 2$):}

\vspace{4pt}
\textbf{Q:} At what point does the curve cross the $y$-axis
(where $x = 0$)?

Set $x(t) = 0$: \quad $t + 1 = 0 \;\Rightarrow\; t = $ \ans{$-1$}

Find $y$: \quad $y($\ans{$-1$}$) = $ \ans{$-1$}${}^2 - 3 = $ \ans{$-2$}

The curve crosses the $y$-axis at the point \ans{$(0,\,-2)$}.

\vspace{10pt}
\textbf{Guided Practice.}\quad Use $f(t) = (t+1,\; t^2-3)$,
$-2 \le t \le 2$.

\vspace{4pt}
(a)\quad Find all values of $t$ in the domain where $y(t) = -2$.
Show your algebra.

\ans{$t^2 - 3 = -2 \;\Rightarrow\; t^2 = 1 \;\Rightarrow\; t = \pm 1$.
Both $t = -1$ and $t = 1$ are in $[-2,\,2]$, so both are valid.}

(b)\quad What are the coordinates of all points on the curve where
$y = -2$?

\ans{At $t = -1$: $(0,\,-2)$. At $t = 1$: $(2,\,-2)$.}

\end{notesbox}

\vspace{0.08in}

%----------------------------------------------------------------------
% Practice
%----------------------------------------------------------------------
\begin{practicebox}

\textbf{Practice.}\quad Let $g(t) = (2t - 3,\; 4 - t^2)$ for
$-1 \le t \le 2$.

\vspace{4pt}
(a)\quad Complete the table of values.

\vspace{4pt}
\begin{center}
\renewcommand{\arraystretch}{1.6}
\begin{tabular}{|c|c|c|c|}
\hline
$t$ & $x(t) = 2t - 3$ & $y(t) = 4 - t^2$ & $(x,\; y)$ \\
\hline
$-1$ & \ans{$-5$} & \ans{$3$} & \ans{$(-5,\;3)$} \\
\hline
$0$  & \ans{$-3$} & \ans{$4$} & \ans{$(-3,\;4)$} \\
\hline
$1$  & \ans{$-1$} & \ans{$3$} & \ans{$(-1,\;3)$} \\
\hline
$2$  & \ans{$1$}  & \ans{$0$} & \ans{$(1,\;0)$}  \\
\hline
\end{tabular}
\end{center}

\vspace{6pt}
(b)\quad What are the start and end points of the curve?

\ans{Start: $(-5,\,3)$ at $t = -1$. End: $(1,\,0)$ at $t = 2$.}

(c)\quad Solve for $t$ when $x(t) = 1$. State the coordinates of
that point.

\ans{$2t - 3 = 1 \;\Rightarrow\; 2t = 4 \;\Rightarrow\; t = 2$.
The point is $(1,\,0)$ (the end point of the curve).}

\end{practicebox}

\end{document}
